% Source: Solutions/Kanishk-PS2-Q6-8.tex (LaTeX source).
\submittedby{Kanishk}

\paragraph{(1)}
We prove this by induction on $k$.

For $k=1$, the result is clear since $A_1$ is countable.

Suppose that $A_1\times A_2\times\cdots\times A_k$ is countable. Since $A_{k+1}$ is also
countable and the product of two countable sets is countable, we get that
\[
  (A_1\times\cdots\times A_k)\times A_{k+1}
\]
is countable. Therefore, by induction, $\prod_{i=1}^{k}A_i$ is countable for every
$k\in\N$.

\paragraph{(2)}
No, the infinite Cartesian product is not necessarily countable. Consider
\[
  A_n=\{0,1\},\qquad n\in\N.
\]
Suppose that $\prod_{n=1}^{\infty}A_n$ is countable. Then we can list all its elements as
\[
  x_1,x_2,x_3,\ldots,
  \qquad\text{where}\qquad
  x_i=(x_{i1},x_{i2},\ldots).
\]
Define a new sequence $y=(y_1,y_2,\ldots)$ by
\[
  y_n=1-x_{nn}.
\]
Clearly, $y$ belongs to the infinite product, since every entry is either $0$ or $1$. But $y$
is different from every $x_n$, since their $n$-th entries are different.

This is a contradiction. Hence, $\prod_{n=1}^{\infty}A_n$ is not necessarily countable.
